\documentclass[11pt]{amsart}
\usepackage{amsmath,amssymb,amsthm}
\usepackage[margin=1.1in]{geometry}
\usepackage{xurl}
\usepackage{hyperref}

\newtheorem{theorem}{Theorem}[section]
\newtheorem{proposition}[theorem]{Proposition}
\newtheorem{lemma}[theorem]{Lemma}
\newtheorem{corollary}[theorem]{Corollary}
\theoremstyle{remark}
\newtheorem{remark}[theorem]{Remark}

\newcommand{\Pt}{\mathbb{P}_t}
\newcommand{\Et}{\mathbb{E}_t}
\newcommand{\Dscr}{\Delta}

\title[Forests of complete graphs]{Negative edge correlation for weighted spanning forests of complete graphs}
\author{Yitzchak Shmalo}
\date{}

\begin{document}

\begin{abstract}
Let $\Pt$ be the arboreal gas on the complete graph $K_n$: the probability measure on
acyclic edge sets $F$ proportional to $t^{|F|}$. We prove that for every $n$ and every
$t>0$ any two distinct edges $e,f$ satisfy $\Pt(e,f\in F)<\Pt(e\in F)\,\Pt(f\in F)$.
The case $t=1$ is due to Gupta. For two edges sharing a vertex the inequality follows
from the real-rootedness of the degree polynomial of a vertex, which is the
characteristic polynomial of a rank-one perturbation of a Hermite Jacobi matrix. For
disjoint edges the edge marginal satisfies an explicit Riccati equation, and the
covariance is controlled by a barrier argument along it.
\end{abstract}

\maketitle

\section{Introduction}

Let $G=(V,E)$ be a finite graph and $t>0$. The \emph{arboreal gas} (or weighted
spanning forest measure) on $G$ is the probability measure
\[
\Pt(F)=\frac{t^{|F|}}{Z_G(t)},\qquad F\subseteq E \text{ acyclic},
\]
where $Z_G(t)=\sum_F t^{|F|}$. It is the $q\to0$ limit of the random-cluster model with
edge weight $p/(1-p)=qt$. Kahn \cite{Kahn}, Grimmett and Winkler \cite{GW} and
Pemantle \cite{Pemantle} asked whether the uniform spanning forest ($t=1$) has
negatively correlated edges, as uniform spanning trees do \cite{BP,FM}; for the arboreal gas at every
$t>0$ this is Conjecture~1.8 of Bauerschmidt, Crawford, Helmuth and Swan \cite{BCHS}.
The conjecture is open for general graphs; for series--parallel graphs the uniform spanning forest satisfies the stronger Rayleigh property \cite{Erickson}.

On complete graphs several cases are known. Gupta \cite{Gupta} proved negative
correlation of distinct edges of the uniform random forest of $K_n$ for every $n$, and Cao and Fu \cite{CF} did the same for the complete bipartite graphs $K_{m,n}$. Tang and Zhang \cite{TZ2} proved it on $K_n$ for $n$ large for uniform spanning forests with a fixed number of components.
Huang \cite{Huang} proved it for adjacent edges at sufficiently large $t$ on any fixed
graph, and on $K_n$ for $n$ sufficiently large when $t$ is sufficiently small or
sufficiently large. \c{C}i\c{c}eksiz and Ravichandran \cite{CR} proved negative edge
correlation for the arboreal gas on any finite graph at sufficiently large fugacity,
under a nondegeneracy condition, and for the random-cluster model on $K_n$ with
$0<q<1$ and common edge weight at least one; the arboreal gas is the limit in which the
weight tends to zero, so it lies outside that range. Tang and Zhang \cite{TZ} prove
negative edge correlation for the random-cluster model with $0\le q<1$ and fixed edge
weight on sufficiently large regular graphs of diverging degree satisfying an
isoperimetric condition. We prove the complete-graph case at
every fugacity.

\begin{theorem}\label{thm:main}
Let $n\ge3$ and $t>0$. For any two distinct edges $e,f$ of $K_n$,
\[
\Pt(e\in F,\ f\in F)<\Pt(e\in F)\,\Pt(f\in F).
\]
\end{theorem}

By symmetry there are two cases: $e$ and $f$ share a vertex (Section~\ref{sec:adj}),
or they are disjoint, which requires $n\ge4$ (Sections~\ref{sec:counts}--\ref{sec:barrier}).
For disjoint edges the proof has three ingredients: exact formulas for the four forest
sums in terms of one sequence of Hermite-type polynomials, a Riccati equation for the
edge marginal $\rho(t)=\Pt(e\in F)$, and the observation that the covariance is a fixed
quadratic $F_t$ evaluated at $\rho(t)$. Positivity of $F_t(\rho(t))$ is then propagated
from both endpoints $t\to0$ and $t\to\infty$ by a first-zero argument, which needs the
sign of the derivative of $F_t(\rho)$ along each root of $F_t$. These signs come from
two polynomial identities and the location of the zeros of a discriminant and of a
cubic.

\section{Generating functions}\label{sec:gf}

Throughout, $a=1/t$. Let $T(w)$ be the exponential generating function of rooted
labelled trees and $U(w)$ that of unrooted labelled trees (both without the empty tree).
Then $T=we^{T}$, $wU'=T$ and $U=T-T^2/2$, since both sides of the last identity vanish
at $0$ and have derivative $T/w$ (using $(1-T)T'=T/w$).

For $b,a\in\mathbb{R}$ define polynomials $P_k=P_k(b,a)$ by
\[
e^{bu-au^2/2}=\sum_{k\ge0}P_k\,\frac{u^k}{k!}.
\]
Differentiating in $u$ gives $P_0=1$, $P_1=b$ and
\begin{equation}\label{eq:rec}
P_k=bP_{k-1}-a(k-1)P_{k-2}\qquad(k\ge2).
\end{equation}
Put $Q_N=P_N-NP_{N-1}$.

\begin{lemma}\label{lem:subst}
For every power series $g$ and every $N\ge0$,
\[
[w^N]\,e^{aU(w)+zT(w)}g(T(w))=[u^N]\,(1-u)\,g(u)\,e^{(N+a+z)u-au^2/2}.
\]
\end{lemma}

\begin{proof}
Write the coefficient as a residue in $w$ and substitute $w=ue^{-u}$, so that $T=u$,
$U=u-u^2/2$ and $dw/w^{N+1}=(1-u)e^{Nu}\,du/u^{N+1}$. This is the Lagrange
inversion formula as a formal identity.
\end{proof}

\subsection*{The degree polynomial of a vertex}
Fix a vertex $v$ of $K_n$, put $N=n-1$, and let
$H(z)=\sum_F t^{|F|}z^{\deg_F(v)}$. Removing $v$ from a forest leaves a forest on the
other $N$ vertices in which every component is joined to $v$ by at most one edge. A
component on $k$ vertices that is not joined to $v$ has weight $t^{k-1}$; a component
joined to $v$ through a chosen vertex has weight $t^{k-1}\cdot tz=zt^k$. By the
exponential formula and Lemma~\ref{lem:subst} with $g=1$,
\begin{equation}\label{eq:H}
H(z)=N!\,t^N[w^N]\,e^{aU(w)+zT(w)}=t^N\,Q_N(N+a+z,\,a).
\end{equation}
In particular $Z:=Z_{K_n}(t)=H(1)=t^{n-1}Q_{n-1}(n+a,a)$.

\section{Edges with a common vertex}\label{sec:adj}

\begin{lemma}\label{lem:real}
All zeros of $H$ are real and negative, and they are not all equal when $N\ge2$.
\end{lemma}

\begin{proof}
Let $J_N$ be the $N\times N$ symmetric tridiagonal matrix with zero diagonal and
off-diagonal entries $\sqrt a,\sqrt{2a},\dots,\sqrt{(N-1)a}$, and let $e_N$ be the last
coordinate vector. Expanding along the last row and using \eqref{eq:rec} gives
$\det(bI-J_N)=P_N(b,a)$. The determinant of $bI-J_N-Ne_Ne_N^{\top}$ is affine in its
last diagonal entry, with cofactor $\det(bI-J_{N-1})=P_{N-1}$, so
\[
\det(bI-J_N-Ne_Ne_N^{\top})=P_N(b,a)-NP_{N-1}(b,a)=Q_N(b,a).
\]
Hence by \eqref{eq:H} the zeros of $H$ are $\lambda-N-a$ for the eigenvalues $\lambda$
of the real symmetric matrix $J_N+Ne_Ne_N^{\top}$, and are real. The coefficients of $H$
are nonnegative and $H(0)\ge1$, so $H$ has no zero in $[0,\infty)$. For $N\ge2$ this
matrix has a nonzero off-diagonal entry, so it is not scalar and its eigenvalues are not
all equal.
\end{proof}

\begin{proposition}\label{prop:adj}
For $n\ge3$, $t>0$ and distinct edges $e,f$ with a common vertex,
$\Pt(e,f\in F)<\Pt(e\in F)\Pt(f\in F)$.
\end{proposition}

\begin{proof}
The leading coefficient of $H$ is $t^N$ (the star at $v$), so by Lemma~\ref{lem:real}
$H(z)=t^N\prod_{j=1}^N(z+\lambda_j)$ with $\lambda_j>0$ not all equal. Let
$Y=\deg_F(v)$ and $p_j=1/(1+\lambda_j)$. Logarithmic differentiation at $z=1$ gives
\[
\mu:=\Et Y=\sum_jp_j,\qquad \Et[Y(Y-1)]=\frac{H''(1)}{H(1)}=\mu^2-\sum_jp_j^2 .
\]
Since the $p_j$ are not all equal, $\sum_jp_j^2>\mu^2/N$, and so
$\Et[Y(Y-1)]<(N-1)\mu^2/N$. The $N$ edges at $v$ are exchangeable; with
$\rho=\Pt(e\in F)$ and $\eta=\Pt(e,f\in F)$ we have $\mu=N\rho$ and
$\Et[Y(Y-1)]=N(N-1)\eta$. Hence $\eta<\rho^2$.
\end{proof}

\section{Forest sums with prescribed edges}\label{sec:counts}

From now on $n\ge4$, $a=1/t$, $b=n+a$, and all $P_k$ are evaluated at $(b,a)$. Let $E$,
$L$ and $D$ be the sums of $t^{|F|}$ over forests containing a fixed edge, two fixed
edges with a common vertex, and two fixed disjoint edges, respectively, and put
$\rho=E/Z$ and $h=(n-2)(n-3)$.

\begin{lemma}\label{lem:counts}
\[
Z=t^{n-1}Q_{n-1},\quad E=t^{n-1}Q_{n-2},\quad L=t^{n-1}Q_{n-3},\quad
D=t^{n-1}\bigl(P_{n-3}-(n-4)P_{n-4}\bigr).
\]
\end{lemma}

\begin{proof}
The formula for $Z$ is \eqref{eq:H} at $z=1$. Contract a prescribed edge $uv$ to a
vertex $m$. A forest containing $uv$ corresponds to a forest of the contracted graph in
which each edge from $m$ to an ordinary vertex $x$ stands for one of the two edges
$ux,vx$; at most one of them can be present. So $m$ is a vertex of degree fugacity
$z=2$ among $n-2$ ordinary vertices, and \eqref{eq:H} with $N=n-2$, $z=2$, times the
weight $t$ of $uv$, gives $E$. Contracting two prescribed edges with a common vertex
gives a vertex of fugacity $3$ among $n-3$ ordinary vertices, and $L$ in the same way.

For two disjoint prescribed edges, contract each to a vertex, $m_1$ and $m_2$, each of
fugacity $2$, among $N=n-4$ ordinary vertices; the edge $m_1m_2$ stands for four edges
and has weight $4t$. If $m_1,m_2$ lie in different components, the components are
counted by $e^{aU+4T}$. If they lie in one component, the path between them has
$k\ge0$ ordinary internal vertices; its edges have total weight $4t^{k+1}$ and each
internal vertex with the trees hanging from it contributes $T/t$, giving $4tT^k$. Hence
the contracted sum is $N!\,t^N[w^N]e^{aU+4T}\bigl(1+4t/(1-T)\bigr)$, and by
Lemma~\ref{lem:subst},
\[
D=t^{n-2}\bigl((1+4t)P_{n-4}-(n-4)P_{n-5}\bigr)=t^{n-1}\bigl(P_{n-3}-(n-4)P_{n-4}\bigr),
\]
the second form by \eqref{eq:rec} with $k=n-3$. For $n=4$ the term with $P_{-1}$ is
absent and both forms equal $t^2(1+4t)$.
\end{proof}

Write $P=P_{n-1}$ and $R=P_{n-2}$. Eliminating $P_{n-3},P_{n-4}$ with \eqref{eq:rec}
expresses the four quantities of Lemma~\ref{lem:counts}, divided by $t^{n-1}$, as
\[
P-(n-1)R,\quad \frac{P-nR}{a},\quad \frac{nP-(n^2+2a)R}{a^2(n-2)},\quad
\frac{(n-4)(n-2)aR+(a-n(n-4))(bR-P)}{a^2h}.
\]
Collecting the coefficients of $P$ and $R$ then verifies the identities
\begin{align}
\frac LZ&=\frac{(2+nt)\rho-2t}{n-2},\label{eq:L}\\
\frac DZ&=\frac{1-2(n-4)t+\bigl(-1/t+2n-9+n(n-4)t\bigr)\rho}{h}.\label{eq:D}
\end{align}

\begin{lemma}[Riccati equation]\label{lem:ric}
For $n\ge4$ and $t>0$,
\[
2t^2\rho'=t-2nt^2+B\rho-n(n-1)t\rho^2,\qquad B=n^2t^2+(2n+1)t-1 .
\]
\end{lemma}

\begin{proof}
Let $M=\binom n2$. Since $t\,\partial_t$ applied to $\sum_Ft^{|F|}$ counts edges,
exchangeability gives $tZ'=ME$. Applied to $E$ it counts, besides the prescribed edge,
the $2(n-2)$ edges sharing a vertex with it and the $h/2$ disjoint ones, so
$tE'-E=2(n-2)L+hD/2$. Hence
$t\rho'=\rho+2(n-2)L/Z+hD/(2Z)-M\rho^2$, and \eqref{eq:L}--\eqref{eq:D} give the claim.
\end{proof}

Define
\[
\delta=th,\quad \alpha=t-2(n-4)t^2,\quad \beta_0=-1+(2n-9)t+n(n-4)t^2,\quad
F_t(r)=\delta r^2-\beta_0r-\alpha .
\]
By \eqref{eq:D}, $F_t(\rho)=\delta(\rho^2-D/Z)$. So, writing $f(t)=F_t(\rho(t))$, the
disjoint case of Theorem~\ref{thm:main} is the statement that $f>0$ on $(0,\infty)$.

\section{The two ends}\label{sec:ends}

\begin{lemma}\label{lem:ends}
For each $n\ge4$, $f(t)/t^5\to2h$ as $t\to0^+$ and $f(t)\to4h/n^4$ as $t\to\infty$.
\end{lemma}

\begin{proof}
Write $Z=\sum Z_kt^k$ and similarly for $E,D$. Counting edge sets with $k$ edges and
removing those containing a cycle,
\[
Z_0=1,\ Z_1=M,\ Z_2=\tbinom M2;\quad E_1=1,\ E_2=M-1,\ E_3=\tbinom{M-1}2-(n-2);
\]
\[
D_2=1,\ D_3=M-2,\ D_4=\tbinom{M-2}2-2(n-2)-2,
\]
with $E_0=D_0=D_1=0$. For $E_3$ the excluded sets are the $n-2$ triangles through the
prescribed edge; for $D_4$ they are the $2(n-2)$ triangles through one prescribed edge
completed by the other, and the two $4$-cycles through both. Multiplying out,
$E^2-ZD=2t^4+O(t^5)$, and $\rho=t+O(t^2)$, so $f=th(E^2-ZD)/Z^2=2ht^5+O(t^6)$.

As $t\to\infty$, $\rho$ is a rational function of $t$ with $\rho\to2/n$, the edge
density of a uniform spanning tree. Write $\rho=2/n+r_1/t+r_2/t^2+O(t^{-3})$. The
coefficients of $t$ and of $t^0$ in Lemma~\ref{lem:ric} give $1+6/n+n^2r_1=0$ and
$-2r_1=n^2r_2+(5-2n)r_1-2/n$, so $r_1=-(n+6)/n^3$ and $r_2=(42-5n)/n^5$. Inserting these
into \eqref{eq:D} gives $\rho^2-D/Z=4/(n^4t)+O(t^{-2})$, hence $f\to4h/n^4$.
\end{proof}

\section{Polynomial facts}\label{sec:poly}

Let $n\ge7$ and
\[
\varphi(t,r)=t-2nt^2+Br-n(n-1)tr^2,\qquad G(t,r)=2t^2\partial_tF_t(r)+\partial_rF_t(r)\,\varphi(t,r).
\]
By Lemma~\ref{lem:ric}, $f'(t)=G(t,\rho(t))/(2t^2)$. Put $A=n(n-4)$,
\begin{align*}
\mathcal L&=n^3(n-4)t^4+2n(n^2-4n-6)t^3+3(n^2-9n+3)t^2+2(n^2-9n+9)t-(2n-3),\\
\mathcal C&=2n^2(n-4)t^3+3n(n-6)t^2+(2n^2-16n+15)t-(2n-3),\\
J&=2A^2t^3-6A(n-2)t^2+6(n-1)(n-3)t-(2n-3).
\end{align*}

\begin{lemma}\label{lem:ids}
As polynomials in $n,t,r$,
\begin{align}
hG&=F_t(r)\bigl[-2hn(n-1)tr+2ht+2hB-n(n-1)\beta_0\bigr]+2(\mathcal Lr-t\mathcal C),\label{eq:id1}\\
\delta t^2\mathcal C^2-\beta_0t\mathcal L\mathcal C-\alpha\mathcal L^2&=-h^2t^5J.\label{eq:id2}
\end{align}
\end{lemma}

\begin{proof}
Both are verified by expanding the definitions; the accompanying scripts do so in exact
arithmetic.
\end{proof}

Let $\ell=2\mathcal L/h$ and $c_0=-2t\mathcal C/h$. By \eqref{eq:id1}, $G(t,r)=\ell r+c_0$
at every real root $r$ of $F_t$. Let $\Dscr=\beta_0^2+4\delta\alpha$ be the
discriminant of $F_t$. When $\Dscr\ge0$ write $r_-\le r_+$ for the roots and
$g_\pm(t)=G(t,r_\pm(t))$. Since $r_-+r_+=\beta_0/\delta$ and $r_-r_+=-\alpha/\delta$,
identity \eqref{eq:id2} gives
\begin{equation}\label{eq:prod}
g_-g_+=-\frac{4t^4J}{h}.
\end{equation}
The leading coefficient of $\ell$ in $t$ is $2n^3(n-4)/h>0$.

\begin{lemma}\label{lem:J}
For $n\ge7$, $J$ has exactly one positive zero $\beta$. It satisfies $\beta>1/(n-3)$,
$J<0$ on $(0,\beta)$, $J>0$ on $(\beta,\infty)$ and $J'(\beta)>0$.
\end{lemma}

\begin{proof}
One checks $J'(t)=6\bigl[(At-(n-2))^2-1\bigr]$, with positive zeros $t_a=(n-3)/A$ and
$t_b=(n-1)/A$, and $J(t_a)=-(n-6)/(n-4)<0$, $J(t_b)=-(n-2)/n<0$, $J(0)=-(2n-3)<0$. So
$J$ is negative on $[0,t_b]$ and strictly increasing on $(t_b,\infty)$ to $+\infty$.
Finally $t_b>1/(n-3)$ because $(n-1)(n-3)-n(n-4)=3$.
\end{proof}

\begin{lemma}\label{lem:disc}
For $n\ge7$, let $t_c=1/(n-3)$. Then $\Dscr$ has exactly two positive zeros
$\tau_1<t_c<\tau_2$, both simple; $\Dscr>0$ on $(0,\tau_1)\cup(\tau_2,\infty)$,
$\Dscr<0$ on $(\tau_1,\tau_2)$, $\Dscr'(\tau_1)<0$ and $\Dscr'(\tau_2)>0$.
\end{lemma}

\begin{proof}
Expanding,
\[
\Dscr=1+(18-4n)t+(6n^2-48n+105)t^2+(-4n^3+38n^2-136n+192)t^3+n^2(n-4)^2t^4 .
\]
Let $c=n-3\ge4$. Under $t=u/(c(1+u))$, which maps $(0,\infty)$ onto $(0,t_c)$, the
polynomial $c^4(1+u)^4\Dscr$ has coefficients, in increasing powers of $u$,
\[
c^4,\ 6c^3,\ c^2(6c+15),\ -2c(c+1)(2c-9),\ -3(c-3)(c+1).
\]
The sequence has exactly one sign change, so by Descartes' rule there is exactly one
positive zero, and it is simple. Under $t=(1+u)/c$, which maps $(0,\infty)$ onto
$(t_c,\infty)$, the coefficients of $c^4\Dscr$ are
\[
-3(c-3)(c+1),\ 2(2c^3-13c^2+3c+18),\ 9(2c^3-5c^2-2c+6),\ 6(c-1)(3c^2-c-6),\ (c+3)^2(c-1)^2,
\]
and with $c=x+4$ the third and fourth are positive, since
$2c^3-5c^2-2c+6=2x^3+19x^2+54x+46$ and $3c^2-c-6=3x^2+23x+38$. Again there is exactly
one sign change, and one simple positive zero. At $t=t_c$, $\Dscr=-3(c-3)(c+1)/c^4<0$.
The signs at $0$, at $t_c$ and at $\infty$ give the rest.
\end{proof}

\section{The barrier argument}\label{sec:barrier}

\begin{proposition}\label{prop:dis}
For $n\ge4$, $t>0$ and disjoint edges $e,f$, $\Pt(e,f\in F)<\Pt(e\in F)\Pt(f\in F)$.
\end{proposition}

\begin{proof}
For $n=4,5,6$, Lemma~\ref{lem:counts} gives
\begin{align*}
E^2-ZD&=t^4(2+4t),\\
E^2-ZD&=t^4(2+24t+84t^2+100t^3),\\
E^2-ZD&=t^4(2+48t+462t^2+2172t^3+5184t^4+5184t^5),
\end{align*}
which are positive for $t>0$. Let $n\ge7$; we show $f>0$ on $(0,\infty)$.

\emph{Signs at the double roots.} Write
$F_t(r)=\delta\bigl(r-\beta_0/(2\delta)\bigr)^2-\Dscr/(4\delta)$. At a zero $\tau$ of
$\Dscr$ the double root is $r^*=\beta_0/(2\delta)$, where $\partial_rF=0$ and
\[
G(\tau,r^*)=2\tau^2\partial_tF_\tau(r^*)=-\frac{\tau^2\Dscr'(\tau)}{2\delta(\tau)} .
\]
By Lemma~\ref{lem:disc} this is positive at $\tau_1$ and negative at $\tau_2$. At
$\tau_2$, \eqref{eq:prod} gives $0<G(\tau_2,r^*)^2=-4\tau_2^4J(\tau_2)/h$, so
$J(\tau_2)<0$ and $\tau_2<\beta$ by Lemma~\ref{lem:J}.

\emph{The interval $(0,\tau_1]$.} Here $t<t_c<\beta$, so $g_-g_+>0$ by \eqref{eq:prod};
neither factor vanishes, and both tend to the positive value $G(\tau_1,r^*)$ at
$\tau_1$. Hence $g_\pm>0$ on $(0,\tau_1]$. By Lemma~\ref{lem:ends}, $f>0$ near $0$. If
$f$ had a first zero $t_0\le\tau_1$, then $f'(t_0)\le0$; but $\rho(t_0)$ is a root of
$F_{t_0}$, so $f'(t_0)=g_\pm(t_0)/(2t_0^2)>0$. Thus $f>0$ on $(0,\tau_1]$.

\emph{The interval $(\tau_1,\tau_2)$.} Here $\Dscr<0$ and $\delta>0$, so $F_t>0$ on all
of $\mathbb R$, and $f>0$.

\emph{The lower root above $\tau_2$.} We show $g_-<0$ on $[\tau_2,\infty)$. On
$(\tau_2,\beta)$, $g_-g_+>0$, and both factors are negative at $\tau_2$, so
$g_\pm<0$ on $[\tau_2,\beta)$. For $t>\beta$, $g_-g_+<0$, and by \eqref{eq:id1}
\begin{equation}\label{eq:diff}
g_+-g_-=\ell\,(r_+-r_-).
\end{equation}
If $\ell(t)=0$ for some $t>\beta$ then $g_+=g_-$ and $g_-g_+\ge0$, which is false; as
$\ell$ is eventually positive, $\ell>0$ on $(\beta,\infty)$. At $\beta$, $\Dscr>0$ since
$\beta>\tau_2$, so $g_\pm$ are differentiable near $\beta$, and $g_-g_+$ has a simple
zero there with derivative $-4\beta^4J'(\beta)/h<0$. If $\ell(\beta)=0$ then
$g_-(\beta)=g_+(\beta)=0$ and the derivative of $g_-g_+$ at $\beta$ would vanish; so
$\ell(\beta)>0$. By continuity $g_\pm(\beta)\le0$, their product is $0$ and their
difference is positive, so $g_-(\beta)<0=g_+(\beta)$. For $t>\beta$, \eqref{eq:diff}
and $g_-g_+<0$ give $g_-<0<g_+$.

\emph{The interval $[\tau_2,\infty)$.} Since $\rho\le2/n$ (a forest has at most $n-1$
edges) while $r_+(t)\to\infty$, and $f>0$ for large $t$ by Lemma~\ref{lem:ends}, we have
$\rho<r_-$ for all large $t$. Suppose $f$ vanishes somewhere in $[\tau_2,\infty)$, and
let $t_0$ be its largest zero there. On $(t_0,\infty)$ the roots are real and distinct
and $f\ne0$, so $\rho<r_-$ throughout, and $\rho(t_0)=r_-(t_0)$. Positivity of $f$ to
the right of $t_0$ requires $f'(t_0)\ge0$, while $f'(t_0)=g_-(t_0)/(2t_0^2)<0$. Hence
$f>0$ on $[\tau_2,\infty)$.
\end{proof}

Theorem~\ref{thm:main} follows from Propositions~\ref{prop:adj} and~\ref{prop:dis}; for
$n=3$ every pair of edges has a common vertex.

\begin{remark}
The argument uses only the exact counts of Lemma~\ref{lem:counts} and properties of
explicit polynomials. Whether the same method applies to other vertex-transitive
graphs, or to unequal edge weights, is not addressed here.
\end{remark}

\section*{Acknowledgments and funding}

This research was supported by the European Research Council (ERC) under the European Union's Horizon Europe research and innovation programme (grant agreement No.~101041711), by the Simons Foundation as part of the Collaboration on the Mathematical and Scientific Foundations of Deep Learning, by Heights Labs, by Convex Nexus Capital, by the Israel Science Foundation (grant number 2258/19), and by the Israel Science Foundation (ISF Grant 4101/25).

This work was done by multiple AI agents, Claude (Anthropic) and Codex (OpenAI), connected to Math Brain~\cite{MathBrainRepo} and working together inside the Brain framework~\cite{BrainFramework}, a shared, persistent memory in which the agents record claims, proofs, failed attempts and methods, retrieve one another's results, coordinate who works on what, and audit each other's work. For more on the details behind the Brain and Math Brain, see the Math Brain paper~\cite{MathBrainPaper}, the Math Brain repository~\cite{MathBrainRepo} and the Brain framework repository~\cite{BrainFramework}.

\section*{Declaration of competing interest}

The author holds equity in Heights Labs and in Convex Nexus Capital and has a management role at Convex Nexus Capital. These entities are included in the funding declaration above.

\section*{Declaration of generative AI and AI-assisted technologies}

OpenAI Codex systems developed the mathematical argument and the exact verification code. Anthropic Claude AI systems audited the proofs, rechecked the algebra with independent code, compared the result with the literature, and wrote and prepared this manuscript and its publication files. The author selected the problem and directed the work; the author did not carry out the mathematics or the writing. Responsibility for publication rests with the author. The AI-assisted internal reviews and compilation checks are not external peer review or formal verification.

\section*{Data and code availability}

The scripts that verify Lemmas~\ref{lem:counts}--\ref{lem:disc} and the base cases, and
that check Theorem~\ref{thm:main} by exhaustive enumeration for $n\le7$, are in the
repository
\url{https://github.com/yspennstate/kn-forest-negative-correlation}.

The paper is part of the Math Brain project, whose public repository is \url{https://github.com/yspennstate/math-brain-public}~\cite{MathBrainRepo}.

\begin{thebibliography}{9}

\bibitem{BCHS}
R.~Bauerschmidt, N.~Crawford, T.~Helmuth and A.~Swan, Random spanning forests and
hyperbolic symmetry, \emph{Comm. Math. Phys.} 381 (2021), 1223--1261.

\bibitem{BP}
R.~Burton and R.~Pemantle, Local characteristics, entropy and limit theorems for
spanning trees and domino tilings via transfer-impedances, \emph{Ann. Probab.} 21
(1993), 1329--1371.
\bibitem{CF}
Y.~Cao and H.~Fu, Negative edge correlation in complete bipartite graphs,
arXiv:2610.36749, 2026.
\bibitem{CR}
R.~A.~\c{C}i\c{c}eksiz and M.~Ravichandran, Negative correlation for forests and the
$q<1$ random-cluster model, arXiv:2610.00549, 2026.

\bibitem{Erickson}
A.~Erickson, Sums of squares and negative correlation for spanning forests of series
parallel graphs, arXiv:1008.3660, 2010.
\bibitem{FM}
T.~Feder and M.~Mihail, Balanced matroids, in \emph{Proceedings of the 24th Annual ACM
Symposium on Theory of Computing}, ACM, New York, 1992, 26--38.
\bibitem{GW}
G.~R.~Grimmett and S.~N.~Winkler, Negative association in uniform forests and connected
graphs, \emph{Random Structures Algorithms} 24 (2004), 444--460.

\bibitem{Gupta}
A.~Gupta, Edges of the uniform random forest of $K_n$ are pairwise negatively
correlated for every $n$, arXiv:2609.26134, 2026.

\bibitem{Huang}
X.~Huang, On negative correlation of arboreal gas on some graphs, arXiv:2311.00965, 2023.

\bibitem{Kahn}
J.~Kahn, A normal law for matchings, \emph{Combinatorica} 20 (2000), 339--391.

\bibitem{Pemantle}
R.~Pemantle, Towards a theory of negative dependence, \emph{J. Math. Phys.} 41 (2000),
1371--1390.

\bibitem{TZ}
P.~Tang and Z.~Zhang, Negative correlation for the random-cluster model below one on
high-degree regular graphs, arXiv:2609.11578, 2026.

\bibitem{MathBrainRepo} Y.~Shmalo, \emph{Math Brain: public snapshot, code and tests}, 2026, \url{https://github.com/yspennstate/math-brain-public}.

\bibitem{MathBrainPaper} Y.~Shmalo, \emph{Math Brain: shared, auditable memory for multi-agent mathematical research}, October 2026, \url{https://github.com/yspennstate/math-brain-public/tree/main/papers/math-brain-2026-10-24}.

\bibitem{BrainFramework} Y.~Shmalo, \emph{Brain Framework}, public software repository, 2026, \url{https://github.com/yspennstate/brain-framework-public}.

\bibitem{TZ2}
P.~Tang and Z.~Zhang, Pairwise negative correlation for uniform spanning subgraphs of
the complete graph, arXiv:2603.10738, 2026.
\end{thebibliography}

\end{document}
